
% ============================================================================
\clearpage
\section*{v4.8.0-r12 Source-Only: Pre-Registered Neutrino Tests, the Octonionic Flavour Group, and Reclassification of the PMNS Claims}
\addcontentsline{toc}{section}{v4.8.0-r12 Source-Only: Pre-Registered Neutrino Tests, Octonionic Flavour Group, PMNS Reclassification}
\label{sec:v480-r12-neutrino-flavour}
% ============================================================================

\paragraph{Purpose.}
Layer r11 introduced pre-registration.  This layer records the seven further pre-registered tests
H2--H8, which apply the triole form to neutrinos and to lepton mixing, the structural facts they rest on,
and a reclassification of the earlier PMNS statements of the Companion.  For every test the hypothesis,
data and decision criterion were committed to the repository before the computation; the commits are
listed below.  Earlier layers are not rewritten.

\noindent\textbf{Source-only status.}
This v4.8.0-r12 layer is source-only.  It does not build a PDF, does not assign a new DOI,
and does not authorize a public release.  The v4.8.0 Version DOI
\href{https://doi.org/10.5281/zenodo.23365244}{10.5281/zenodo.23365244} and the Concept DOI
\href{https://doi.org/10.5281/zenodo.19210728}{10.5281/zenodo.19210728} are unchanged by it.

\subsection*{v4.8.0-r12 structural facts}
\label{subsec:v480-r12-structure}

\paragraph{Theoremlet v4.8.0-r12.1 (triole spectrum and sign).}
For \(v_k=t\bigl(1+\sqrt2\,r\cos(\delta+2\pi k/3)\bigr)\) one has
\(v_0v_1v_2=t^3\bigl(1-\tfrac32r^2+\tfrac{r^3}{\sqrt2}\cos3\delta\bigr)\).  For \(r=1\) at most one
component is negative, and exactly one is negative if and only if \(\cos3\delta<1/\sqrt2\); the boundary
\(\delta=\pi/12\) is the point where one mass vanishes.  The vector \(v\) is the spectrum of the Hermitian
circulant matrix \(t\,\mathbf 1+\tfrac{c}{\sqrt2}P+\tfrac{\bar c}{\sqrt2}P^2\) (\(P\) the cyclic shift),
whose eigenvectors are the fixed Fourier vectors.

\paragraph{Proof sketch.}
Expand the product with \(\sum_k\cos\theta_k=0\), \(\sum_k\cos^2\theta_k=3/2\) and
\(\prod_k\cos\theta_k=\tfrac14\cos3\delta\) for \(\theta_k=\delta+2\pi k/3\).  The circulant statement is the
discrete Fourier diagonalization.  The identities were also checked numerically to \(10^{-15}\).  \(\square\)

\paragraph{Theoremlet v4.8.0-r12.2 (octonionic flavour group).}
In the Cayley--Dickson convention of layer r04 the automorphisms of \(\OO\) that permute
\(e_1,\dots,e_7\) up to sign form a group of order \(1344\).  Its stabilizer of \(+e_7\) has order \(96\),
acts unitarily with determinant one and irreducibly on \(\mathbb C^3=\mathrm{span}_{\mathbb C}(e_4,e_5,e_6)\)
with complex structure \(L_{e_7}\), and has the same element-order statistics and the same number of
conjugacy classes (ten) as \(\Delta(96)\subset SU(3)\).

\paragraph{Proof sketch.}
An automorphism is fixed by the images of \(e_1,e_2,e_4\); all \(14\cdot12\cdot8\) choices extend and
preserve the full multiplication table.  Restriction, unitarity, determinant, irreducibility
(\(\langle|\chi|^2\rangle=1\)) and the class data were computed exhaustively (scripts
\texttt{k22\_oktonion\_triplett.py}, \texttt{k23\_test\_h6.py}).  The identification with \(\Delta(96)\) is by
these invariants; an explicit isomorphism is not written out here.  \(\square\)

\paragraph{Theoremlet v4.8.0-r12.3 (residual reflections).}
With the triole \(e_4\to e_5\to e_6\) as charged-lepton symmetry and the electron along its invariant
direction, the fifteen involutions of the stabilizer preserve the following columns of the lepton mixing
matrix:
\begin{center}
\small
\begin{tabular}{>{\raggedright\arraybackslash}p{0.40\textwidth}>{\raggedright\arraybackslash}p{0.07\textwidth}>{\raggedright\arraybackslash}p{0.20\textwidth}>{\raggedright\arraybackslash}p{0.21\textwidth}}
\toprule
Involution & count & preserved \(|U_{\alpha i}|^2\) & pattern\\
\midrule
pure sign changes (one generation fixed) & 3 & \((\tfrac13,\tfrac13,\tfrac13)\) & TM2\\
Fano centre in \(\{e_1,e_2,e_3\}\), first lift & 3 & \((\tfrac23,\tfrac16,\tfrac16)\) & TM1\\
Fano centre in \(\{e_1,e_2,e_3\}\), second lift & 3 & \((0,\tfrac12,\tfrac12)\) & \(\theta_{13}=0\)\\
Fano centre in \(\{e_4,e_5,e_6\}\) & 6 & \((\tfrac13,\tfrac{2\mp\sqrt3}{6})\) & \(\Delta(96)\) column\\
Fano centre \(e_7\) & 0 & --- & only order-4 lifts\\
\bottomrule
\end{tabular}
\end{center}
The TM1 lift and the \(\theta_{13}=0\) lift over the same Fano image differ by a pure sign change.

\paragraph{Theoremlet v4.8.0-r12.4 (Cayley--Dickson involution as CP).}
The involution \(p+q\ell\mapsto p-q\ell\) commutes with every lift of the triole and with all TM1 and
\(\theta_{13}=0\) involutions, and acts on the triplet as \(z\mapsto-\bar z\) in the basis in which the
triole is the plain cyclic shift.  If the Majorana mass matrix is invariant under it and under a residual
reflection, then \(\sin^2\theta_{23}=\tfrac12\) and \(\delta_{CP}\in\{90^\circ,270^\circ\}\) exactly.

\paragraph{Proof sketch.}
\(X=-\mathbf1\) forces a real mass matrix in the generation basis; the Fourier matrix of the triole turns
this into the \(\mu\)--\(\tau\) reflection, whose consequences are classical.  Checked on 2000 random
points of the four-dimensional solution space for each reflection type (script \texttt{k26\_test\_h8.py}).
\(\square\)

\subsection*{v4.8.0-r12 pre-registered tests}
\label{subsec:v480-r12-tests}

Data: NuFIT 6.0 (IC24 with SK atmospheric data, normal ordering unless stated), DESI DR2 with CMB for
\(\sum m_\nu<0.0642\) eV, KamLAND-Zen for \(m_{\beta\beta}\); Gaussian approximation with asymmetric
errors where a \(\chi^2\) is quoted.
\begin{center}
\small
\begin{tabular}{>{\raggedright\arraybackslash}p{0.06\textwidth}>{\raggedright\arraybackslash}p{0.10\textwidth}>{\raggedright\arraybackslash}p{0.36\textwidth}>{\raggedright\arraybackslash}p{0.38\textwidth}}
\toprule
Test & commit & hypothesis & result under the fixed criterion\\
\midrule
H2 & \texttt{434c5a8} & neutrinos are an \(r=1\) triole, signs free & compatible: normal ordering, \(\sum m_\nu=0.0592\pm0.0002\) eV; inverted ordering disfavoured (\(0.102\) eV); weakly discriminating\\
H3 & \texttt{e3b140e} & triole signs are Majorana CP parities & compatible, \(m_{\beta\beta}=2.7\pm0.7\) meV (\(\le3.5\) meV); not testable in the foreseeable future\\
H4 & \texttt{00d297e} & TM1 against TM2, equal prior & TM1 strongly preferred, \(\chi^2=4.4\) vs.\ \(14.7\) (2 d.o.f.); the \(\theta_{12}\) part was not blind\\
H5 & \texttt{16eff15} & polar angle of the preserved column is a Hopf angle of the Companion & failed: \(0/32\) canonical combinations compatible; free azimuth degenerate\\
H6 & \texttt{e1e2393} & list of algebraic selection criteria for the TM1 lift & main criterion (dihedral \(S_3\) with the triole) selects the \(\theta_{13}=0\) lift, i.e.\ speaks against TM1\\
H7 & \texttt{e6da7db} & Heisenberg compensator \(J\) as sign automorphism selects the lift & no selection for any \(J\); the candidates negating \(e_1\) and fixing \(e_7\) do not commute with the triole\\
H8 & \texttt{f6734ab} & Cayley--Dickson involution as CP symmetry & tension, \(\chi^2=8.8\) (3 d.o.f.); fixed predictions \(\sin^2\theta_{23}=\tfrac12\), \(\delta_{CP}=270^\circ\), with TM1 \(\sin^2\theta_{12}=0.318\)\\
\bottomrule
\end{tabular}
\end{center}

\paragraph{Reading.}
The octonionic structure already present in the Companion produces the known mixing patterns TM1, TM2
and \(\Delta(96)\) as candidates and, through the Cayley--Dickson involution, a definite CP structure.  It
does not select TM1: neither minimality, nor the cascade projections (they commute with all sign changes),
nor the compensator distinguishes the TM1 lift from the excluded one.  The predictions of H4 and H8 will be
tested by JUNO (\(\theta_{12}\)) and by DUNE and Hyper-Kamiokande (\(\theta_{23}\), \(\delta_{CP}\)).

\subsection*{v4.8.0-r12 reclassification of earlier PMNS statements}
\label{subsec:v480-r12-reclassification}

\paragraph{Proposition~\ref{prop:pmns-z3}.}
The value \(\arctan\sqrt2=54.74^\circ\) is not the maximal mixing angle; maximal mixing is \(45^\circ\).
The shifted value \(\arctan(\sqrt2)\,\delta_{\rm CL}=52.1^\circ\) gives \(\sin^2\theta_{23}=0.622\), against
\(0.470\) (normal ordering, \(\theta_{23}=43.3^\circ\)) about \(9\sigma\) in the Gaussian approximation and
about \(6\sigma\) against the inverted-ordering value \(0.550\).  The proposition and the
\(\theta_{23}^{\rm PMNS}\) line of the unified prediction table are reclassified as \emph{refuted}.  The
replacement statement of this layer is Theoremlet~r12.4 with the open selection problem.

\paragraph{Theorem~\ref{thm:op1op4-unified}.}
Its CKM and PMNS lines rest on the continuum-limit conjecture for \(\delta_{\rm CL}\), whose \(\beta\)-formula
was classified in layer r09 as a coincidence with an exponent chosen after the fact.  The CKM line inherits
that status; the PMNS line is refuted as above.  The CKM phase remark (\(\pi/3\) against \(\approx70^\circ\))
is unaffected and remains open.

\subsection*{v4.8.0-r12 ledger and guard update}
\label{subsec:v480-r12-ledger-guard}

\[
R_{\rm mix}^{v480\text{-}r12}=(0,0,1)
\]
for (protocols committed before computation; all seven hypotheses decided; selection principle for the
residual reflection still open).

The layer spends no guard:
\[
R_{\rm glob}^{(v480\text{-}r12,+)}=(0,1,1,1,1),
\qquad
\Gamma_{\rm guard}^{v480\text{-}r12}=(1,1,1,1,1),
\qquad
\beta_4=1/256\ \text{unspent}.
\]
It does not imply exact finite \(C5\) zero closure, exact all-readout terminality,
\(C6/T2\), a global photon-clock/RH theorem, or global all-readout closure.
